\scanpage{26}
\uuline{changement de base} $f^*$, \uuline{et compatible} % L1
\uuline{également avec} $\mathbb R f_*$ [N.B. $T_\ell$ s'étend % L2
% slanted
\margin{compatibilité avec $f^!$, $\mathbb R f_!$, avec l'hom résidu etc...} % L2a
évidemment en un foncteur % L3
$D^b(\mathcal M(X))\unsure{-|\to}D^b(\mathcal M_\ell(X))$]. $\{$ La détermination % L4
des $T_\ell$ est encore ramenée au cas où % L5
$X$ est de t.f. sur $\mathbb Z$. % L6

% slanted
\margin{\uline{NB.} Ceci exclut déjà le choix trivial $\mathcal M^+(X)=0$ pour tt $X$, car il faudrait qu'on ait $\mathbb R^{\unsure{*|x}}f_*(\mathbb Z_{\underline{\ell}})=0$ pour $f,\ell$, ce qui n'est pas vrai en général ---} % L7
\deleted{Si $X$ est un $\mathbb Q$-schéma, \inserted{de \deleted{\unsure{présentation}} \inserted{type} fini,} alors on % L7a
a aussi un foncteur % L8
\[
T_\infty=T_\infty^{(X)}:D^b(\mathcal M(X))\longrightarrow\mathcal M_\infty(X) % L9
\]
$\mathcal M_\infty(X)$ est la catégorie des \inserted{$\mathscr O_X$-}Modules % L10
cohérents. Ce foncteur $T_\infty$ est encore % L11
\uline{exact} et \uline{fidèle}, compatible avec % L12
$f^*$ et $\mathrm Rf_*$.} % L13
\uline{Signaler aussi la compatibilité de} $T_\ell$ \uline{avec l'isom. de Künneth} % L13a

5) Pour tout $X$, on a un \deleted{\unsure{un}} % L14
élément canonique % L15
\[
\deleted{\mathbb Z}\mathbb Q_X(-1)\ \text{ou}\ \mathbb 1_X(-1)\ \in\mathrm{Ob}\,\mathcal M^+(X) % L16
\]

% archive sheet number: 25
